\section{Отчёт о выполнении задачи 1}
\label{sec:task1}

В рамках выполнения первой задачи было проведено изучение теоретических основ эффективности
критериев согласия по Ходжесу--Леману.

Были выполнены следующие действия:
\begin{itemize}
    \item Сопоставлены эффективность по Питмену, по Бахадуру и по Ходжесу--Леману.
    Обоснован выбор именно эффективности по Ходжесу--Леману для данной работы.
    \item Определены необходимые величины и допущения.
    \item Проанализированы особенности оценки малых вероятностей. Изучена проблема
    того, что при больших $n$ величина $1-\pi_n(\theta)$ становится малой и при
    конечном числе моделирований может оказаться равной $0$. Определены подходы к
    выбору «линейного участка» зависимости для корректной оценки наклона.
    \item Выполнен обзор литературы и инструментов.
\end{itemize}

\subsection{Полученные артефакты}

\begin{itemize}
  \item Репозиторий отчётности.
  \item Перечень необходимых данных для эксперимента. Определён минимальный набор величин, который потребуется при реализации эксперимента в \texttt{PySATL-experiment}.
  \item Список литературы для дальнейшего изучения. Сформирован приоритетный список источников, необходимый для углублённого разбора при реализации.
\end{itemize}

\subsection{Ожидаемые артефакты}

\begin{itemize}
  \item Схема архитектуры предлагаемого решения.
  \item Pull request в \texttt{pysatl-experiment} с устранением ограничений.
\end{itemize}

\clearpage
